Let \(\mathscr K\) be the locally uniform closure appearing in \(\text{def:contour-bank-handle}\).  We first verify its compactness and that the defining restrictions survive closure.  Uniform sub-Gaussianity gives constants \(c,C>0\), depending only on \(\psi_{\eta}\), such that every underlying law satisfies
\[
 P(|\eta|>u)\le C e^{-cu^2/\psi_{\eta}^{2}}.
\]
Thus the laws are uniformly tight.  From any sequence one obtains a weakly convergent subsequence by diagonal extraction of the distribution functions at rational continuity points.  The same tail bound implies uniform integrability of \(e^{S|\eta|}\) for each fixed \(S<\infty\).  Truncation at \(|\eta|\le L\), followed by weak convergence and then \(L\to\infty\), shows pointwise convergence of the MGFs on \(\mathbb C\).  The derivative envelope
\[
 \sup E_P[|\eta|e^{S|\eta|}]<\infty
\]
makes the transforms equicontinuous on \(\{|z|\le S\}\); hence the convergence is locally uniform.  The same uniform-integrability argument applies to every moment, so centering and \(|\kappa_k|\ge\delta\) pass to the limit.  The \(\psi_2\) bound passes by Fatou's inequality.  Therefore \(\mathscr K\) is sequentially compact in the metrizable compact-open topology and hence compact.

Fix \(M\in\mathscr K\).  The zero-localization lemma gives a zero in \(\{|z|\le R_0\}\).  Since \(M(0)=1\), \(M\) is not identically zero, and its zeros are isolated.  We may therefore choose a rational \(\rho(M)\in(R_0,R_1)\) whose circle contains no zero.  Set
\[
 m(M)=\inf_{|z|=\rho(M)}|M(z)|>0.
\]
The circle encloses at least the localized zero.  Factoring the finitely many enclosed zeros with their multiplicities and applying the Cauchy coefficient identity to the remaining zero-free analytic factor gives
\[
 \frac{1}{2\pi i}\oint_{|z|=\rho(M)}\frac{M'(z)}{M(z)}\,dz
 =N_{|z|=\rho(M)}(M)\ge1.
\]

Consider the compact-open neighborhood
\[
 \mathscr U_M
 =\left\{\widetilde M\in\mathscr K:
   \sup_{|z|=\rho(M)}|\widetilde M(z)-M(z)|<m(M)/2\right\}.
\]
Every \(\widetilde M\in\mathscr U_M\) has boundary modulus at least \(m(M)/2\).  Moreover the homotopy \(M_t=(1-t)M+t\widetilde M\) is zero-free on that circle for every \(t\in[0,1]\).  The integral
\[
 t\longmapsto \frac{1}{2\pi i}\oint_{|z|=\rho(M)}\frac{M_t'(z)}{M_t(z)}\,dz
\]
is continuous and integer-valued, so it is constant.  Hence every member of \(\mathscr U_M\) has the same positive enclosed zero count as \(M\).

The neighborhoods \(\{\mathscr U_M:M\in\mathscr K\}\) cover the compact set \(\mathscr K\).  Choose a finite subcover \(\mathscr U_{M_1},\ldots,\mathscr U_{M_J}\), put
\[
 \rho_j=\rho(M_j),
 \qquad
 a_{\star}=\frac12\min_{1\le j\le J}m(M_j)>0,
\]
and let \(C_j=\{|z|=\rho_j\}\).  Any admissible treatment transform belongs to some \(\mathscr U_{M_j}\), and for that \(j\),
\[
 N_{C_j}(M)\ge1,
 \qquad
 \inf_{z\in C_j}|M(z)|\ge a_{\star}.
\]
Only the fixed class constants enter \(\mathscr K\), so \(J\), the radii, and \(a_{\star}\) do not depend on \(n\).